% !TEX program = pdflatex
% ============================================================
% MCM/ICM English paper template.
% Compile: pdflatex main.tex (twice)  or  compile.py main.tex
% ============================================================
\documentclass[12pt]{article}

\usepackage[a4paper,margin=1in]{geometry}
\usepackage{amsmath,amssymb,amsthm,bm,amsfonts}
\usepackage{graphicx}
\usepackage{booktabs}
\usepackage{longtable}
\usepackage{multirow}
\usepackage{caption}
\usepackage{subcaption}
\usepackage{float}
\usepackage{enumitem}
\usepackage{listings}
\usepackage{xcolor}
\usepackage{algorithm}
\usepackage{algpseudocode}
\usepackage{fancyhdr}
\usepackage{hyperref}
\hypersetup{colorlinks=true,linkcolor=black,citecolor=blue,urlcolor=blue}

% ====== 页面约束宏包 ======
\usepackage{flafter}
\usepackage{afterpage}
\usepackage{placeins}
\usepackage{needspace}
\usepackage{tocloft}

% ====== 核心页面控制 ======
\raggedbottom
\setlength{\parskip}{0.5ex plus 0.2ex minus 0.2ex}
\setlength{\abovecaptionskip}{6pt}
\setlength{\belowcaptionskip}{4pt}

% 浮动体间距紧凑化
\setlength{\floatsep}{8pt plus 2pt minus 2pt}
\setlength{\textfloatsep}{10pt plus 2pt minus 3pt}
\setlength{\intextsep}{8pt plus 2pt minus 2pt}
\setlength{\dblfloatsep}{8pt plus 2pt minus 2pt}
\setlength{\dbltextfloatsep}{10pt plus 2pt minus 3pt}

% 浮动体参数
\setcounter{topnumber}{3}
\setcounter{bottomnumber}{2}
\setcounter{totalnumber}{5}
\renewcommand{\topfraction}{0.85}
\renewcommand{\bottomfraction}{0.6}
\renewcommand{\textfraction}{0.15}
\renewcommand{\floatpagefraction}{0.7}

% 防止孤行
\widowpenalty=10000
\clubpenalty=10000

% 容忍度
\tolerance=800
\emergencystretch=1.5em
\hbadness=3000

% ---------- 目录紧凑化 ----------
\renewcommand{\cftdot}{}
\setlength{\cftbeforesecskip}{2pt}
\setlength{\cftbeforesubsecskip}{1pt}
\renewcommand{\cftsecfont}{\normalsize}
\renewcommand{\cftsubsecfont}{\small}
\renewcommand{\cftsecpagefont}{\normalsize}
\renewcommand{\cftsubsecpagefont}{\small}

% ---------- 图表与代码 ----------
\graphicspath{{figures/}{../figures/}}
\captionsetup{font=small}
\lstset{basicstyle=\ttfamily\small,numbers=left,numberstyle=\tiny\color{gray},
  frame=single,breaklines=true,showstringspaces=false,tabsize=4}

\newtheorem{theorem}{Theorem}
\newtheorem{assumption}{Assumption}

% ---- Control number header (fill in) ----
\pagestyle{fancy}\fancyhf{}
\rhead{Team \# XXXXXXX}\lhead{Problem X}\cfoot{\thepage}

\begin{document}

% ===================================================================
%   Summary Sheet  （⛔ HARD: must fit on ONE page）
%   If it overflows: compress text, reduce to \small, or cut verbose sentences.
% ===================================================================
\begin{center}
  {\Large\textbf{Title of Your Paper}}\\[0.5em]
  \textbf{Summary}
\end{center}

\small  % Use small font to fit more content on one page
\noindent
State the problem, your modeling approach, key methods, and the most important
quantitative results and conclusions here. Judges read this first --- make every
sentence count and include concrete numbers.

\bigskip
\noindent\textbf{Keywords:} keyword1; keyword2; keyword3; keyword4

% Summary Sheet done — force new page
\newpage

% ===================================================================
%   Table of Contents  （⛔ HARD: must fit on ONE page）
%   If too long: \setcounter{tocdepth}{1} (sections only) or {0}
% ===================================================================
\setcounter{tocdepth}{2}       % Default to subsections; reduce if overflows
\tableofcontents
\newpage

% ------------------ 正文 ------------------
\section{Introduction}
\subsection{Background}
\subsection{Restatement of the Problem}
\subsection{Our Work}
\FloatBarrier

\section{Assumptions and Justifications}

\subsection{Core Assumptions}
Each assumption must state its \textbf{rationale} and discuss its \textbf{limitation}
--- what happens if relaxed, and by how much the conclusion would change.

\begin{assumption}
\textbf{Assumption 1 (\ldots\ assumption).} Rationale: \ldots.
Limitation: if this does not hold (e.g., \ldots), the model would
\ldots\ (overestimate / underestimate / break down) in \ldots\ scenarios.
A possible remedy is \ldots.
\end{assumption}

\begin{assumption}
\textbf{Assumption 2 (\ldots\ assumption).} Rationale: \ldots. Limitation: \ldots.
\end{assumption}

\subsection{Sensitivity of Assumptions}
\begin{enumerate}[label=(\arabic*)]
  \item \textbf{If Assumption 1 fails}: \ldots\ (quantify bias magnitude, direction,
        whether the conclusion reverses).
  \item \textbf{If Assumption 2 fails}: \ldots
\end{enumerate}
\textbf{Conclusion}: The critical assumption is \ldots; if it fails, conclusions
are no longer reliable. Relaxing other assumptions has limited impact
(within $\pm\cdots\%$).
\FloatBarrier

\section{Notation}
\begin{center}
\begin{tabular}{cl}
\toprule
Symbol & Description \\
\midrule
$x_i$ & \ldots \\
\bottomrule
\end{tabular}
\end{center}
\FloatBarrier

\section{Model Development}
\subsection{Model I}
\begin{equation}
\min_x f(x)=\sum_i c_i x_i \quad \text{s.t.}\ \sum_i a_{ij}x_i\le b_j .
\end{equation}

% ---- Figures: always use [htbp] not [H] to avoid large gaps ----
\begin{figure}[htbp]
\centering
\includegraphics[width=0.75\textwidth]{fig_convergence.png}
\caption{Algorithm convergence curve.}
\label{fig:conv}
\end{figure}

As shown in Figure~\ref{fig:conv}, the algorithm converges to the optimum
within approximately …… iterations, with stable convergence behavior.
\FloatBarrier

% ============================================================
\section{Model Comparison and Validation}
% ⛔ Key scoring point: compare against 1-2 baselines/ablations
% ============================================================

\subsection{Baseline Models}
To validate the effectiveness of our model, we compare against:
\begin{enumerate}[label=(\arabic*)]
  \item \textbf{Our Model}: \ldots\ (brief description).
  \item \textbf{Baseline A (\ldots)}: \ldots\ (classical method / simpler alternative).
  \item \textbf{Ablation B (w/o \ldots\ module)}: \ldots\ (tests a specific component).
\end{enumerate}

\subsection{Multi-Metric Comparison}

\begin{table}[htbp]
\centering
\caption{Multi-metric comparison across models.}
\label{tab:compare}
\begin{tabular}{lcccc}
\toprule
Model & Metric 1 (\ldots) & Metric 2 (\ldots) & Metric 3 (\ldots) & Overall \\
\midrule
Our Model    & …… & …… & …… & \textbf{……} \\
Baseline A   & …… & …… & …… & …… \\
Ablation B   & …… & …… & …… & …… \\
\bottomrule
\end{tabular}
\end{table}

From Table~\ref{tab:compare}:
\begin{enumerate}[label=(\arabic*)]
  \item Our model outperforms Baseline A on Metric 1 by \textbf{……\%},
        because \ldots.
  \item Removing the \ldots\ module (Ablation B) drops Metric 2 by \textbf{……\%},
        confirming its critical contribution.
  \item All models perform similarly on Metric 3, suggesting it is insensitive
        to model choice.
\end{enumerate}

% Radar chart for visual comparison
\begin{figure}[htbp]
\centering
\includegraphics[width=0.6\textwidth]{fig_radar_compare.png}
\caption{Multi-metric radar comparison across models.}
\label{fig:radar_compare}
\end{figure}
\FloatBarrier

\section{Sensitivity Analysis}

\begin{figure}[htbp]
\centering
\includegraphics[width=0.75\textwidth]{fig_sensitivity.png}
\caption{Parameter sensitivity analysis (tornado chart).}
\label{fig:sens}
\end{figure}

As seen in Figure~\ref{fig:sens}, parameter …… has the greatest influence
on the objective. When perturbed by $\pm 20\%$, the output varies by ……,
while other parameters show minimal impact, confirming model robustness.
\FloatBarrier

% ============================================================
\section{Strengths, Weaknesses, and Future Work}
% ============================================================

\subsection{Strengths}
\begin{enumerate}[label=(\arabic*)]
  \item \textbf{Theoretical novelty}: \ldots\ (e.g., combining …… with ……).
  \item \textbf{Methodological innovation}: \ldots\ (e.g., novel algorithm / mechanism).
  \item \textbf{Practical utility}: \ldots\ (e.g., low computational cost, accessible inputs).
  \item \textbf{Interpretability}: \ldots\ (e.g., each parameter has clear meaning).
\end{enumerate}

\subsection{Weaknesses and Improvement Directions (In-Depth)}
% ⛔ Don't just write "parameter sensitivity" or "data limitation" --- be specific
\begin{enumerate}[label=(\arabic*)]
  \item \textbf{Assumption limitation}: Assumption …… may fail in \ldots\ scenarios,
        causing errors up to ……\%. Improvement: introduce …… correction /
        nonlinear extension / robust optimization.
  \item \textbf{Data dependence}: The model is sensitive to …… data (per sensitivity
        analysis). Current data precision is ……; if error grows to ……,
        conclusions may reverse. Improvement: Bayesian inference /
        interval estimation instead of point estimates.
  \item \textbf{Computational bottleneck}: When problem scales to ……,
        current $O(\cdots)$ complexity becomes prohibitive.
        Improvement: …… (approximation / parallelization / heuristic).
  \item \textbf{Generalization limit}: The model assumes …… domain; applying to
        …… requires recalibrating …… parameters.
\end{enumerate}

\subsection{Future Extensions}
\needspace{3\baselineskip}
…… (e.g., apply to …… scenarios, extend to multi-stage decisions, fuse with ……)

\section{Conclusion}
……

% ------------------ 参考文献 ------------------
\begin{thebibliography}{99}
\bibitem{ref1} Author. Title. Publisher/Journal, Year.
\end{thebibliography}

% ------------------ 附录 ------------------
\appendix
\section{Code}
\lstinputlisting[language=Python]{../code/solve_q1.py}

\end{document}
